\chapter*{前言}
\addcontentsline{toc}{chapter}{前言}

\begin{notebox}
\paragraph{从一次水位观测开始}

一条库水位观测从测站出发，经过采集网关、通信网络和数据服务，最终成为监测页面上的一个点。值班员点击这个点时，应能查到测点位置、观测时间、质量码（标记这条观测可信与否的代码）和历史变化；水位超过阈值时，系统还要生成预警，记录复核与处置过程。

完成这样的功能，需要把水文与工程知识写成数据字段、计算规则和程序流程。本书面向已学过一门编程语言、尚未系统学习软件工程与前后端开发的水利专业学生，从一次数据请求入手，在提供的工程骨架上逐步实现页面、接口、数据库与三维场景之间的协作。
\end{notebox}

\section*{课程定位}

《国家水网建设规划纲要》提出“建设数字孪生水网”\cite{nationalWaterNetwork2023}。在工程实践中，监测感知、数据治理、模型计算与运行管理需要共同支撑防洪、供水和工程安全等任务。平台开发因此涉及两类相互关联的工作：解释水利业务的依据与约束，并将其实现为能够运行和检验的软件。

本课程衔接程序设计基础与水利工程应用。全书围绕一个虚构的案例水库安全监测平台展开：28个测点的观测数据怎样从接口到页面、到曲线、到三维场景，再到预警和处置。学完核心内容后，你应当能够：

\begin{enumerate}
\item 说清一个智慧水利平台由哪些部分组成，一次“查询渗压计 PZ-07 最新观测”的请求经过了哪些环节；
\item 在配套工程骨架上实现监测对象查询、观测曲线、三维测点关联和一条简单的预警规则；
\item 解释水利观测数据的单位、时间、质量码、坐标与高程在程序里怎样表示，为什么不能省；
\item 遇到“页面空白”“接口 404”“曲线断线”这类常见异常时，知道先看什么证据。
\end{enumerate}

课程会让你接触 Vue、Spring Boot、Three.js、PostgreSQL 等多种技术，但目标是\textbf{读得懂、改得动、讲得清}，不是熟练掌握整套技术栈。消息队列、数字孪生调度、复杂性能调优等内容安排为选读。

\section*{内容与学习顺序}

全书九章，按照\textbf{“打基础 $\rightarrow$ 学技术 $\rightarrow$ 做项目”}的思路组织，分为三个部分。

\begin{tipbox}
\paragraph{第一部分\quad 基础篇（第1--3章）}

第1章说明智慧水利平台由什么组成，并不写代码地追踪一次查询请求；第2章用案例灌区练习需求获取、结构化分析和需求规格说明；第3章为案例水库划分模块、约定接口，并沿接口走查两条业务路径。这部分产出的需求表、模块表和接口契约是后面各章实现的依据。
\end{tipbox}

\begin{tipbox}
\paragraph{第二部分\quad 技术篇（第4--7章）}

四章各解决一个问题，每章都从一个能运行、能修改的最小程序开始：
\begin{itemize}
\item \textbf{前端}（第4章）：HTML、CSS 与 JavaScript 基础，一个读取接口并显示四种状态（加载、有数据、空数据、错误，以及切换对象后的过期响应）的监测页面，再用 Vue 3 组织路由与共享状态
\item \textbf{后端}（第5章）：第一个返回固定数据的 HTTP 接口，接入 PostgreSQL 数据库，拆成控制器、服务与数据访问三层，再加上 JWT 认证与基本授权；消息队列的完整实现为选学
\item \textbf{三维场景}（第6章）：用几何体搭出坝体，加载现成模型，处理坐标与高程，把28个测点绑定到场景对象上
\item \textbf{观测数据展示}（第7章）：画出第一条真实观测曲线，处理缺测与质量码，实现曲线与三维测点的双向联动；降采样与性能调优为进阶
\end{itemize}
\end{tipbox}

\begin{tipbox}
\paragraph{第三部分\quad 应用篇（第8--9章）}

第8章用前面各章做出的模块完成一条业务链：渗压计 PZ-07 的一次异常观测经过质量检查、预警定级、值班员确认和工单处置，最后能回查当时的依据；部署与冒烟测试作为指导实践。8.5节的数字孪生案例是独立选读。第9章按课程目标回顾各章成果，并给出判断一项新技术是否值得投入的方法。
\end{tipbox}

\section*{怎样读这本书}

\paragraph{阅读规则}每章的每一节开头有一行“本节层次”，把小节分成三类。\textbf{核心}：能解释原理、完成小规模实现和基本题，是课堂讲授与闭卷考核的范围。\textbf{指导实践}：在工程骨架上修改、运行和验证，主要在实验课完成。\textbf{拓展}：分析适用条件与代价，供课程设计与自学。阅读规则只有一条：按章顺序读各章的核心小节和指导实践，标为拓展的小节在完成对应阶段之后再回来读。核心小节不依赖拓展小节；拓展小节常常依赖核心小节的运行结果。

\paragraph{阶段与版本}第4--7章开头各有一段“工程版本线”，第8章开头的“配套工程入口”说明 v4 到 v5，把案例平台从空骨架 v0 做到完整平台 v5；配套工程按阶段包 S0--S6 组织，每个阶段都有一个能运行的入口和一份验收要求。表\ref{tab:preface-version-stage}把两套编号对在一起：正文里的“v2”和配套里的“S3”指的是同一件事的两个侧面，前者说平台做到了哪一步，后者说从哪个文件启动、验收什么。核心路线就是沿这张表从 S0 走到 S6。表中两个词先在这里说明：“契约”指前后端事先写下并共同遵守的接口约定（路径、字段、错误体）；“竞态”指连续发出两次请求后，先发出的响应反而后到，页面把旧结果当成新结果显示。

\begin{table}[htbp]
\centering
\small
\caption{工程版本 v0--v5 与配套阶段包 S0--S6 的对应}
\label{tab:preface-version-stage}
\begin{tabular}{p{0.8cm}p{1.6cm}p{2.0cm}p{4.6cm}p{4.2cm}}
\toprule
版本 & 阶段 & 章节 & 配套入口 & 达到的能力 \\
\midrule
v0 & S0 & 第1章 1.2.1 & \texttt{S0-demo-record.md} & 解释一次查询的六个环节；不写代码 \\
v1 & S1 & 4.2--4.4 & \texttt{lesson44.html}、\texttt{lesson44-detail.html} & 静态列表与详情页，筛选与排序，只需浏览器 \\
v1 & S2 & 4.5、4.6.1--4.6.2（核心）；4.6.3--4.7（指导实践） & \texttt{lesson45.html}；\texttt{index.html} 与 \texttt{src/router/} 骨架 & 对着教学接口显示四种页面状态、处理竞态；Vue 路由与登录守卫 \\
v2 & S3 起点 & 5.2 & \texttt{Lesson52Application} & 第一个接口：固定数据、路径与查询参数、契约错误体 \\
v2 & S3 中点 & 5.4.1--5.4.2 & \texttt{Lesson54Application} & 数据来自 PostgreSQL，三层结构；无登录 \\
v2 & S3 终点 & 5.5--5.6、8.3 & 完整 \texttt{backend/} 与 Compose & 认证与授权齐备，承接 S2 页面联调 \\
v3 & S4 & 6.1 & \texttt{lesson61.html} & 坝体几何体与28个测点绑定 \\
v4 & S5 & 7.2--7.4 & \texttt{lesson74.html} & 观测曲线与三维测点双向联动 \\
v5 & S6 & 8.4、8.6 部署 & \texttt{classify.js}、\texttt{closeloop-check.mjs}、\texttt{smoke.sh} & 质量门禁、四级定级、确认与工单闭环；部署冒烟测试 \\
\bottomrule
\end{tabular}
\end{table}

\paragraph{运行、修改与验证}学习代码清单时，先按说明运行并核对输出，再修改一个输入或配置，预测并观察结果，最后注入故障、说明原因。章末思考题用于检查概念之间的联系，实践题要求提交代码与验证记录。记录包括输入、预期结果、实际结果和失败时的排查过程。

\paragraph{几条具体建议}
\begin{enumerate}
\item \textbf{准备环境的时机。}S1、S2、S4、S5 只需要 Node.js 与浏览器；Java 17 与 Maven 从5.2节起需要；Docker（容器运行环境）与数据库从5.4节起需要。不必在开课第一周装好全部环境。
\item \textbf{组队做课程设计。}条件允许时，3--4人一组分工完成一个小型监测平台：前端、后端、三维场景各有分工，在协作中体会接口契约为什么要先于代码写下来。
\item \textbf{查阅文档并验证。}遇到报错时，先记录完整错误和依赖版本，再查相应版本的官方文档。采用社区方案前，检查其输入条件与版本是否相符，并用同一故障重测；在实验记录中保留来源和验证结果。
\end{enumerate}

\section*{推荐教学组织方式}

三个学时方案按学生最终能做到什么来编制，学时少的方案少走几个阶段，已走过的阶段仍然完整。方案之间的差别在于沿表\ref{tab:preface-version-stage}走到哪个阶段：32学时走到 S3 中点，48学时走到 S5，56学时走到 S6 并完成部署验收。每个方案都包含表\ref{tab:preface-core-optional}左列的对应部分，右列的选学专题留给课程设计或自学。

\begin{itemize}
\item \textbf{32学时方案（到 S3 中点）：}学生能解释一次请求的全过程，把一句含糊需求写成验收条件，在教学接口上做出带四种状态的监测页面，写出第一个 HTTP 接口并把它接到数据库。不讲三维场景、曲线联动和预警闭环；第6--8章可作为课程设计的阅读材料。表\ref{tab:preface-32-hours}是章节分配，适合专业方向课或选修课。
\begin{table}[htbp]
\centering
\caption{32学时方案章节分配（到 S3 中点）}
\label{tab:preface-32-hours}
\begin{tabular}{p{2.8cm}p{1.4cm}p{8.2cm}}
\toprule
章节或环节 & 学时 & 课堂重点 \\
\midrule
第1章 & 2 & 平台组成、三类规范、一次查询的六个环节（S0）、预警闭环 \\
第2章 & 4 & 用例与异常流、数据字典、从含糊需求到验收条件 \\
第3章 & 4 & 需求分给模块、接口契约、走查一次查询 \\
第4章 4.1--4.5 & 9 & 静态页面（S1）、请求与四种页面状态、竞态故障（S2 核心） \\
第5章 5.1--5.4 & 11 & 第一个接口（S3 起点）、数据库查询与三层（S3 中点） \\
复习与个人答辩 & 2 & 对着自己的页面和接口解释一次失败 \\
\bottomrule
\end{tabular}
\end{table}

\item \textbf{48学时方案（到 S5）：}在32学时的基础上完成认证与授权（S3 终点）、坝体场景与测点绑定（S4）、曲线与三维联动（S5）。第8章只读8.1--8.2，了解业务链怎样由已学模块接成，不做预警与工单的实现。表\ref{tab:preface-48-hours}给出分配。
\begin{table}[htbp]
\centering
\caption{48学时方案章节分配（到 S5）}
\label{tab:preface-48-hours}
\begin{tabular}{p{2.8cm}p{1.4cm}p{8.2cm}}
\toprule
章节或环节 & 学时 & 课堂重点 \\
\midrule
第1章 & 2 & 平台组成、三类规范、一次查询的六个环节（S0）、预警闭环 \\
第2章 & 4 & 用例与异常流、数据字典、从含糊需求到验收条件 \\
第3章 & 3 & 需求分给模块、接口契约、走查两条路径 \\
第4章 & 10 & S1 静态页面、S2 四种状态与竞态、Vue 路由与登录守卫 \\
第5章 5.1--5.6 & 13 & S3 起点、中点与终点：接口、数据库、三层、认证与授权 \\
第6章 6.1--6.2 & 6 & 几何体坝体、模型加载、坐标与高程、测点绑定（S4） \\
第7章 7.1--7.4 & 6 & 观测记录与质量码、第一条曲线、双向联动（S5） \\
第8章 8.1--8.2 & 2 & 读懂 PZ-07 业务链用到哪些已学模块 \\
复习与项目答辩 & 2 & 演示 S5 页面，解释一次切换测点时的请求过程 \\
\bottomrule
\end{tabular}
\end{table}

\item \textbf{56学时方案（到 S6 与部署）：}讲授第1--8章的核心与指导实践小节，第4、5章各占14学时，其中一半用于随堂实验；第8章完成质量门禁、预警定级、确认与工单闭环，并用 Compose 与冒烟测试做一次部署验收。第6章6.3--6.4、第8章8.5和第9章9.2.1--9.2.2作为选读。表\ref{tab:preface-56-hours}给出分配，讲授与实验之比约为30:26。
\begin{table}[htbp]
\centering
\caption{56学时方案章节分配（到 S6 与部署）}
\label{tab:preface-56-hours}
\begin{tabular}{p{2.8cm}p{1.2cm}p{1.2cm}p{7.0cm}}
\toprule
章节或环节 & 讲授 & 实验 & 课堂重点 \\
\midrule
第1章 & 2 & 0 & 一次请求的完整过程、平台分层、三类规范与业务闭环 \\
第2章 & 3 & 2 & 用例、异常流、数据字典与可验证需求 \\
第3章 & 3 & 2 & 需求分给模块、接口契约、两条走查记录、一份 ADR（架构决策记录） \\
第4章 & 7 & 7 & 页面、数据请求与四种状态、Vue 组件、路由与共享状态 \\
第5章 & 7 & 7 & 第一个 HTTP 接口、数据库与三层、校验与事务、认证与权限 \\
第6章 & 3 & 2 & 几何体坝体、坐标与高程、模型加载与测点绑定 \\
第7章 & 3 & 2 & 观测记录与质量码、第一条曲线、测点状态与联动 \\
第8章 & 2 & 3 & 质量门禁、四级定级、确认与工单闭环、部署冒烟测试 \\
第9章与答辩 & 0 & 1 & 按课程成果回顾、现场小修改与解释 \\
\bottomrule
\end{tabular}
\end{table}

\item \textbf{课程设计方案：}以第8章贯穿案例为中心，要求学生分组完成需求说明、架构设计、前端页面、后端接口和三维场景集成，可与毕业设计或综合实训衔接。
\end{itemize}

\paragraph{课堂主线与选学专题}表\ref{tab:preface-core-optional}列出各部分的课堂主线与选学专题。课堂主线包含核心小节及相应的指导实践；选学专题用于课程设计、毕业设计或自学，不进入闭卷考核。

\begin{table}[htbp]
\centering
\caption{核心内容与选学内容划分}
\label{tab:preface-core-optional}
\begin{tabular}{p{2.2cm}p{5.6cm}p{5.6cm}}
\toprule
部分 & 课堂主线（含指导实践） & 选学（课程设计/进阶） \\
\midrule
方法基础 & 第1--3章：平台组成与一次查询、需求方法与灌区案例、模块划分、接口走查与一份 ADR & 第1章智能边界（1.1.4）、第2章螺旋/Scrum/DevOps、第3章内聚耦合度量、ATAM 演练与复杂度分析、3.5.3--3.5.5 \\
前端 & 第4章：HTML/CSS/JS 基础、四种页面状态与竞态处理、request.js 与登录联调、Vue 组件、路由与 Pinia & 层叠上下文与大屏暗色主题（4.3.7--4.3.8）、4.5.8、事件循环（4.5.9）、Vite 工程化（4.8.3）、焦点管理与虚拟列表、构建发布（附录C） \\
后端 & 第5章：REST、JPA、事务、统一异常、JWT认证与基本授权、提交后事件与异步消息的概念（5.7.1） & Kafka 生产与消费的完整实现、事务发件箱、令牌轮换与撤销、可观测性、Testcontainers \\
三维与展示 & 第6章6.1--6.2、第7章第一条曲线、质量码与联动 & 管线与 LOD（6.1.6）、Cesium 场景（6.2.5）、倾斜摄影与BIM（6.3）、数字孪生架构（6.4）、多源数据与时间对齐（7.1.4--7.1.5）、LTTB 降采样与性能预算（7.2.6--7.2.7）、专题地图与色彩（7.2.8--7.2.9）、聚合与触控（7.3.3、7.4.3） \\
综合案例 & 第8章8.1--8.4 的 PZ-07 业务链：质量检查、定级、确认与工单；8.6部署冒烟测试 & Cesium 扩展（8.2.2）、物理设计与缓存一致性（8.3.7--8.3.8）、模型可信度（8.4.3）、8.5数字孪生案例 \\
展望 & 第9章全书回顾与技术评估方法（9.1、9.2.3） & 9.2.1--9.2.2 延伸专题 \\
\bottomrule
\end{tabular}
\end{table}

\section*{贯穿案例卡}

全书用两个教学案例分工协作。\textbf{水利工程安全监测平台}是贯穿主线，载体是虚构的案例水库：第1章用它建立业务闭环认知，第3章为它完成架构设计，第4--8章沿 v0 到 v5 的版本线把它从空骨架逐步做成完整平台，第9章以它复盘课程成果。\textbf{案例灌区用水计划}只服务于第2章的需求分析教学：配水申报的多方冲突比安全监测更适合练习需求获取与建模；第2章末尾的迁移一节把这套方法映射回水库主线。两个案例共享需求分析的方法，业务对象与流程各有侧重。案例水库的工程参数、测点编码、角色和预警等级统一按8.1节使用，表\ref{tab:preface-case-card}摘列常用参数，便于前置章节查阅：

\begin{table}[htbp]
\centering
\caption{贯穿案例参数快照（权威来源：表8.1）}
\label{tab:preface-case-card}
\begin{tabular}{p{3.2cm}p{10.2cm}}
\toprule
项 & 取值 \\
\midrule
工程规模 & 流域面积\cpBasinArea\,km$^2$，总库容\cpTotalCapacity 亿\,m$^3$，坝高\cpDamHeight\,m（坝顶\cpCrestElev\,m、坝基\cpBaseElev\,m） \\
特征水位 & 死水位\cpDeadLevel\,m；汛限水位\cpFloodLimitLevel\,m；正常蓄水位\cpNormalLevel\,m；设计洪水位\cpDesignFloodLevel\,m；校核洪水位\cpCheckFloodLevel\,m \\
泄洪设施 & \cpGateCount 孔弧形闸门，单孔净宽\cpGateWidth\,m，闸底高程\cpGateSillElev\,m，最大开度\cpGateMaxOpen\,m \\
监测体系 & \cpStationCount 个观测点＝\cpPzCount 渗压＋\cpDispCount 位移＋\cpWlCount 库水位＋\cpRfCount 雨量；\cpGatewayCount 台采集网关；贯穿测点\cpKeyStation \\
预警口径 & 蓝、黄、橙、红四级＋无预警（NONE）；证据不足显示“未评估” \\
角色口径 & 值班员、专业分析员、审批人、运维员（第8章另设只读扩展角色审计员） \\
质量码 & valid、suspect、missing 三值 \\
\bottomrule
\end{tabular}
\end{table}

\section*{配套资源}

配套资源包含可运行工程、模拟数据、参考答案和故障练习。目录与用途见表\ref{tab:preface-companion-resources}，版本与入口见表\ref{tab:preface-companion-versions}。实验前先读对应目录的README，按阶段说明启动服务，并用给定请求检查响应。课程报告注明代码提交号、数据版本、环境配置和运行命令，小组成员据此复现实验。

\begin{table}[htbp]
\centering
\caption{配套资源目录与使用方式}
\label{tab:preface-companion-resources}
\begin{tabular}{p{3.2cm}p{4.0cm}p{5.2cm}}
\toprule
资源位置 & 内容 & 使用方式 \\
\midrule
\texttt{companion/}\allowbreak\texttt{water-platform-}\allowbreak\texttt{demo/} & 前端、后端、教学接口、数据库脚本与容器编排；\texttt{STAGES.md} 按 S0--S6 给出入口与验收 & 按 STAGES.md 选择本章阶段的入口启动，逐阶段替换接口或页面 \\
\texttt{companion/}\allowbreak\texttt{datasets/} & 测点元数据、水位、雨量、渗压和预警样例 & 第7章实践题与第8章验收测试使用，字段和单位随README核对 \\
附录A & 习题参考答案、评分要点与自测提示 & 先独立完成，再按题号检索思路；答案不替代实验记录 \\
Issues反馈渠道 & 勘误、环境问题和可复现实验报告 & 在配套仓库Issues提交章节、版本、日志和最小复现步骤 \\
\bottomrule
\end{tabular}
\end{table}

\paragraph{配套工程的使用。}\texttt{companion/water-platform-demo/}提供登录认证、测点与观测查询、消息消费入库、容器启动和前后端测试，并按阶段包组织练习代码。\texttt{STAGES.md}列出每个阶段的入口文件、运行命令、验收和故障练习；\texttt{MAPPING.md}列出书中清单与仓库文件的对应关系，以及工程版本与阶段的对应；\texttt{teaching-api/}是零依赖的教学接口，只需 Node.js 即可启动，前端各阶段不必先装 Java 和数据库。阶段页提供的是可运行的教学切片；详情路由、GLTF 模型与 GIS 集成等任务由读者在骨架上完成，二者的界线见 STAGES.md。

\paragraph{代码仓库结构。}\texttt{frontend/}的可运行入口是\texttt{package.json}、六个阶段页（\texttt{lesson44.html}、\texttt{lesson44-}\allowbreak\texttt{detail.html}、\texttt{lesson45.html}、\texttt{lesson61.html}、\texttt{lesson74.html}和 Vue 工程的\texttt{index.html}）、\texttt{src/router/}\allowbreak\texttt{index.js}、\texttt{src/stores/}\allowbreak\texttt{monitoring.js}、\texttt{src/utils/}\allowbreak\texttt{request.js}和\texttt{src/components/}\allowbreak\texttt{Monitoring}\allowbreak\texttt{Dashboard.vue}，并配有\texttt{vite.config.js}、\texttt{Dockerfile}和\texttt{nginx.conf}；\texttt{backend/}包含\texttt{pom.xml}、\texttt{src/main/java/}\allowbreak\texttt{edu/example/}\allowbreak\texttt{qingyuan/}下的实体、Repository、服务、控制器、JWT 过滤器和安全配置，以及阶段入口\texttt{edu/example/}\allowbreak\texttt{lesson52/}（S3 起点）与\texttt{lesson54/}（S3 中点）；\texttt{db/}包含\texttt{001\_init.sql}与\texttt{002\_seed.sql}；根目录的\texttt{docker-compose.yml}负责编排前端、后端、PostgreSQL、Redis和Kafka，\texttt{secrets/}只提交\texttt{*.example}模板，\texttt{smoke.sh}是部署后的冒烟测试。

\paragraph{示例数据集。}\texttt{companion/datasets/}（仓库提交\texttt{4fed08e}）由\texttt{generate.py}生成，固定随机种子为\texttt{20260807}。\texttt{stations.csv/json}含28个测点元数据，字段为\texttt{asset\_id}、\texttt{asset\_type}、\texttt{display\_name}、单位、经纬度、\texttt{elevation\_m}、量程、精度、阈值和投运日期；\texttt{water\_level.csv}含3个库水位测点各1,200条、5分钟间隔的观测；\texttt{rainfall.csv}含5个雨量测点各288条，\texttt{piezometer.csv}含12个渗压测点各288条；8个位移测点当前仅提供\texttt{stations}中的元数据与阈值，其观测序列留作第7章练习用\texttt{generate.py}扩展生成。三类已生成观测共同使用\texttt{asset\_id}、\texttt{occurred\_at}、\texttt{version}、\texttt{event\_id}、\texttt{value}、\texttt{unit}、\texttt{quality}和\texttt{source}字段；\texttt{warnings.json}含蓝、黄、橙、红、无预警（NONE）与未评估样例。质量码只取\texttt{valid}、\texttt{suspect}、\texttt{missing}，文件中的缺测、超量程和跳变记录均保留原因。时间采用ISO 8601、时区采用UTC+8；报告中注明数据版本。数据集由脚本生成，不是真实工程数据。

\paragraph{环境版本基线。}按表\ref{tab:preface-companion-versions}准备环境，并记录实际安装版本。仓库快照中前端锁定Vue 3.4.38、Vite 5.4.8、Pinia 2.1.7、Vue Router 4.4.5、ECharts 5.5.1、Three.js 0.160.1和\texttt{@vitejs/plugin-vue} 5.2.3；后端\texttt{pom.xml}锁定Java 17、Spring Boot 3.2.10，应用版本为0.1.0；Compose 使用\texttt{timescale/timescaledb-ha:pg16}（同时含 PostGIS 与 TimescaleDB 扩展）、\texttt{redis:7.2-alpine}和\texttt{apache/kafka:3.7.2}。Python生成脚本只使用标准库，建议Python 3.11或更高版本。更新依赖后重新执行阶段测试；操作系统或CPU架构不同导致的问题，提交时附上版本输出和启动日志。

\begin{table}[htbp]
\centering
\caption{配套资源的实际版本快照与入口}
\label{tab:preface-companion-versions}
\begin{tabular}{p{3.1cm}p{5.1cm}p{3.7cm}}
\toprule
资源 & 仓库中的确切入口与版本 & 用途 \\
\midrule
配套代码 & \texttt{companion/}\allowbreak\texttt{water-platform-}\allowbreak\texttt{demo/}；入口见\texttt{STAGES.md}；Vue 3.4.38、Vite 5.4.8、Spring Boot 3.2.10 & 第4--8章各阶段的运行、联调与容器启动 \\
示例数据 & \texttt{companion/datasets/}；仓库提交 \texttt{4fed08e}；生成器种子 \texttt{20260807} & 第7章实践题和第8章验收测试 \\
参考答案 & 附录A及配套网站的习题参考答案；按章号、题号检索 & 按章节和题号核对评分要点、验收证据与延伸阅读 \\
运行编排 & \texttt{docker-compose.yml}；镜像采用本节环境基线；经前端访问 \texttt{/readyz} 与 \texttt{/healthz} & 启动前后端、数据库、缓存和消息服务 \\
\bottomrule
\end{tabular}
\end{table}

\paragraph{答案与勘误。}附录A与配套网站按章节和题号提供客观题答案、简答题评分要点、实践题验收清单与延伸阅读。先独立完成题目并保存运行日志，再按题号核对：结果是否正确，异常输入是否处理，解释是否与代码行为一致。发现文字、公式、代码环境或字段问题时，可在配套仓库Issues提交书名、章节、页码、资源提交号、环境版本和复现步骤；若涉及数据处理，还应附输入样例与预期输出。

\section*{适用对象}

\paragraph{适用读者}本书以水利工程及相关专业高年级本科生为主要读者，兼顾研究生，以及有志于从事水利信息化、智慧水利平台开发的技术人员。

\paragraph{先修知识}学习本书前，你需要具备以下基础：

\begin{itemize}
\item 至少学过一门编程语言（C、Python、Java 或 C\# 均可）
\item 了解基本的数据结构概念（数组、链表、树等）
\item 对水利工程有初步认识（水力学、水文学等先修课程）
\item 了解基本的数据库概念与SQL语言（推荐但不强制）
\end{itemize}

没有接触过 Java 或 JavaScript 的读者，先读附录B的语言衔接单元，再进入第4、5章。本书也可作为水利行业技术人员了解平台开发技术的参考读物。

\section*{致谢}

感谢北京五一视界数字孪生科技股份有限公司提供《51WIM 数字孪生水利平台 V1.0（DEMO）使用说明手册》参考资料，并授权本书8.5节使用相关产品界面截图。书中引用的研究成果、标准与开源项目见参考文献及相应代码说明。

欢迎读者提交阅读与实验中发现的问题。包含章节位置、输入条件和运行结果的反馈，有助于核对内容并修正示例。

\vspace{2em}
\begin{flushright}
编者\\
2026年9月
\end{flushright}
